% !TeX root = ../tesis.tex


\chapter{Trabajo a futuro}

\label{section:apendix1}



En cuanto a la parte experimental, se realizó la comparación con los resultados obtenidos de et al. quienes realizaron mediciones experimentales de la FSC para una muestra de eritrocitos sanos de 34.4 g/dL de CHCM de un citómetro de flujo, asumiendo incidencia en la cara circular del eritrocito ($\mathbf{k_1}$) en función de una incidencia lateral ($\mathbf{k_2}$) iluminada a $\lambda=632.8$ nm. En la Fig. \ref{subfig:STLFPi2} se muestran la FSC correspondiente a la polarización $\mathbf{k_2}$ en funcion de la $\mathbf{k_1}$, donde cada punto representa la medición correspondiente a un eritrocito y las zonas de color amarillo representan una mayor densidad de eritrocitos. En azul claro se resalta el punto correspondiente al cálculo de CST para el modelo de eritrocito empleado en este trabajo con una concentración de 30.6 g/dL. Se observa que los valores calculados que se integraron en $\theta$ de 2.2 a $8.2^{\circ}$ y en $\varphi$ de modo que $|\sin(\varphi)|\geq 2.2^{\circ}/\theta$, para simular las condiciones experimentales bajo las que se hicieron las mediciones. Se observa que el punto obtenido de manera numérica está dentro de las mediciones experimentales.

%
\begin{figure}[h!]
	\centering
	\includegraphics[width=0.43\textwidth]{4-Resultados/FSCpaper.pdf}\label{subfig:STLFPi2}
	\vspace{-0.1cm}
	\caption{ Eficiencia de extinción $Q_{\text{ext}}$ en función de la longitud de onda $\lambda$ de la onda incidente para un eritrocito (puntos azules), un esferocito (puntos rojos) y un microcito (puntos verdes) orientados con su eje de simetría en el eje $z$ iluminado con una onda en dirección del eje $z$ y polarizada en $x$. La línea continua que une los puntos es una guía visual. Se resaltan en gris las bandas que 
		rodean los mínimos asociados a la banda de Soret y a las bandas Q.}
	\label{fig:STLF}
\end{figure}
%

